\chapter{Implementation Details}

\section{Core Packages and Class Modules}
All classes are in the default package, in the \texttt{src} folder. Chapter 3
lists the classes. The three speed classes differ only in the delivery days
and in the fee rule (Table~\ref{tab:fees}). The base cost is 100.

\begin{table}[h]
\centering
\caption{Delivery days and fee of each speed.}
\label{tab:fees}
\renewcommand{\arraystretch}{1.3}
\begin{tabularx}{\textwidth}{|l|c|X|c|}
\hline
\textbf{Speed} & \textbf{Days} & \textbf{Fee rule} & \textbf{Fee}\\
\hline
Standard & 5 & base cost & 100.0\\
\hline
Express & 2 & base cost + base cost / days & 150.0\\
\hline
Overnight & 1 & base cost + base cost / days + 50 & 250.0\\
\hline
\end{tabularx}
\end{table}

\section{Key Algorithms and Logic Implementation}
\textbf{Creating a parcel.} \texttt{ParcelFactory} compares the speed name and
returns the matching object, or \texttt{null} for an unknown or \texttt{null}
speed. \texttt{Tracker} builds the ID from the highest number in use plus one,
so the first ID is \texttt{ORD1001}. An ID is never used twice, also when the
manifest has a gap.

\textbf{Moving a parcel.} \texttt{Route.getNext()} finds the current hub in the
array and returns the next one. At the last hub it returns \texttt{null}, and
the parcel does not move.

\begin{lstlisting}[style=java]
public double calculateTotalFee() {      // in ExpressParcel
    return getBaseCost() + getBaseCost() / getDeliveryDays();
}
\end{lstlisting}

\section{Exception Handling Mechanisms}
\texttt{InvalidTrackingIdException} extends \texttt{Exception} \cite{ref1}. The method
\texttt{checkId()} throws it if the ID is \texttt{null}, does not start with
\texttt{ORD}, has no digit after \texttt{ORD}, or has a character that is not a
digit. The method
\texttt{findParcel()} throws it if no parcel has that ID. \texttt{Main} catches
it, prints the message, and shows the menu again.

\section{File I/O Implementation}
\texttt{ManifestStore} uses \texttt{FileWriter}, \texttt{PrintWriter}, and
\texttt{Scanner} \cite{ref1}. \texttt{saveManifest()} writes all parcels and replaces the
old file. \texttt{appendLog()} opens the log in append mode and adds one line.
At start, \texttt{loadManifest()} reads the manifest and rebuilds each parcel.
\texttt{loadHistory()} then reads the log and gives each parcel its
milestones. If a file does not exist, the program starts with no data.

Every file and every \texttt{Scanner} is opened with try-with-resources. Java
closes it when the block ends, also after an error. \texttt{Main} opens its
\texttt{Scanner} the same way.
